% !TeX program = lualatex
% !TeX encoding = utf8
% !TeX spellcheck = uk_UA
% !TeX root = ../EMConspect.tex

%=========================================================
\chapter{Випромінювання електромагнітних хвиль}\label{\currfilebase}
\graphicspath{{\currfilebase/Pictures}}
%=========================================================
\epigraph{Поле, що спочиває, належить заряду. Поле, що тікає, --- вже ні.}{Переосмислення ідей сучасної фізики}

Ми розглянули основні властивості електромагнітних хвиль: їхню природу, шкалу частот і закономірності поширення. Проте залишається ключове питання: як саме виникають ці хвилі?

%% --------------------------------------------------------
\section{Історична довідка: досліди Герца}
%% --------------------------------------------------------

\begin{wrapstuff}[type=figure, o, top=3, width=0.5\linewidth]
    \centering
    \includegraphics[width=\linewidth]{Hertz_Dipole}
    \caption{Вібратор Герца.}\label{pic:Hertz_vibrator}
\end{wrapstuff}

У 1888~році Генріх Герц\footnote{Генріх Рудольф Герц (1857--1894) --- німецький фізик.} уперше експериментально отримав електромагнітні хвилі. У його досліді джерелом випромінювання слугував так званий \emphx{вібратор Герца} — прямолінійний провідник із невеликим повітряним проміжком посередині (рис.~\ref{pic:Hertz_vibrator}). Завдяки цьому проміжку двом половинам вібратора можна було надати значних зарядів: коли різниця потенціалів сягала граничного значення, у
повітряному зазорі відбувався пробій (проскакувала іскра), і заряди через іонізоване повітря перетікали з однієї половини вібратора на іншу ---
в такому \emphx{відкритому контурі} виникали електромагнітні коливання. Щоб швидкозмінні струми існували лише у вібраторі й не замикалися через
джерело живлення, між вібратором і джерелом вмикали \emphx{дроселі}. Високочастотні коливання існують, поки іскра замикає проміжок; загасання
ж їх відбувається в основному не за рахунок джоулевих втрат (як у закритому коливальному контурі), а за рахунок \emphx{випромінювання
електромагнітних хвиль}.

Для виявлення хвиль Герц застосовував другий (\emphx{приймальний}) вібратор. Під дією змінного електричного поля хвилі, що надходить від
випромінювача, електрони в приймальному вібраторі здійснюють вимушені коливання --- у ньому збуджується швидкозмінний струм. Якщо розміри
приймального вібратора такі самі, як і в випромінювального, то частоти власних коливань у них збігаються, і вимушені коливання за рахунок
\emphx{резонансу} досягають помітної величини. Ці коливання Герц реєстрував по проскакуванню іскри в мікроскопічному зазорі посередині
приймального вібратора або по свіченню мініатюрної газорозрядної трубки, увімкненої між його половинами.

Герц не лише експериментально довів існування електромагнітних хвиль, а й уперше почав вивчати їхні властивості --- поглинання й заломлення в
різних середовищах, відбивання від металевих поверхонь тощо. На досліді вдалося також виміряти швидкість електромагнітних хвиль, яка виявилась
рівною швидкості світла. Саме цей збіг --- швидкості електромагнітних хвиль зі швидкістю світла, виміряною задовго до їх відкриття, --- став
відправним пунктом для ототожнення світла з електромагнітними хвилями й створення \emphx{електромагнітної теорії світла}.

\begin{Attention}
	Електромагнітна хвиля існує \emphx{без джерел поля} в тому сенсі, що після її випромінювання електромагнітне поле хвилі вже не пов'язане з
	джерелом. Саме цим електромагнітна хвиля відрізняється від статичних електричного й магнітного полів, які не існують у відриві від джерела.
\end{Attention}

%% --------------------------------------------------------
\section{Елементарна модель випромінювання}
%% --------------------------------------------------------

\subsection*{Якісна картина}

Розглянемо електричне поле точкового заряду $q$. У стані спокою це звичайне електростатичне поле з радіальними силовими лініями. Нехай у момент
$t=0$ заряд під дією зовнішньої сили починає рухатися з прискоренням $a$, а через невеликий час $\tau$ дія сили припиняється, після чого заряд
рухається рівномірно зі швидкістю $v_q=a\tau$ (рис.~\ref{pic:VelocityGraph}); нехай $v_q\ll c$, а $\tau$ дуже мале порівняно з часом спостереження
$T$.

\begin{SCfigure}[1.75][h!]\centering
	\localinput{velocity-graph.tikz}
	\caption{Швидкість заряду: короткий імпульс прискорення тривалістю $\tau$, після якого заряд
		рухається рівномірно\label{pic:VelocityGraph}}
\end{SCfigure}

Оскільки збурення поля поширюється не миттєво, а зі швидкістю світла $c$, різні спостерігачі дізнаються про прискорення заряду в різний момент часу.
Нехай пройшов час $T\gg\tau$ (рис.~\ref{pic:ElementaryModel}):

\begin{itemize}
	\item \emphx{Зовнішній спостерігач} (на відстані, більшій за $cT$) іще нічого не знає про рух заряду: сигнал про початок руху до нього не
	      встиг дійти, і він бачить поле нерухомого заряду в його \emphx{початковому} положенні.
	\item \emphx{Внутрішній спостерігач} (на відстані, меншій за $c(T-\tau)$) уже знає, що прискорення закінчилось і заряд рухається рівномірно, ---
	      він бачить поле рівномірно рухомого заряду в його \emphx{поточному} положенні.
	\item У проміжній \emphx{тонкій сферичній оболонці} товщиною $c\tau$ --- \emphx{області випромінювання} --- силові лінії поля мусять
	      \enquote{зшивати} ці два різні поля, утворюючи характерний \emphx{злам}.
\end{itemize}

\begin{figure}[h!]\centering
	\localinput{elementary-model.tikz}
	\caption{Заломлені (\enquote{злами}) силові лінії поля короткочасно прискореного заряду:
		межа між \enquote{старим} і \enquote{новим} полем\label{pic:ElementaryModel}}
\end{figure}

\begin{Attention}
	Саме цей злам силових ліній і є, по суті, електромагнітна хвиля, що \enquote{відірвалась} від заряду й розповсюджується самостійно. Оскільки
	лінії неперервні (теорема Гаусса), вони не можуть обірватись на межі --- і саме необхідність \enquote{зшити} два різні поля без розриву ліній
	породжує додаткову, \emphx{поперечну} складову поля, якої немає ні в чисто електростатичному, ні в чисто \enquote{рівномірно-рухомому} полі
	окремо.
\end{Attention}

\subsection*{Математичне виведення поля випромінювання}

Розглянемо детальніше злам силової лінії, що виходила з заряду під кутом $\theta$ до напряму його прискорення (рис.~\ref{pic:ElementaryDerivation}).
Нехай спостереження ведеться в момент $T=r/c$ (де $r$ --- відстань до точки спостереження), тобто поле, яке зараз \enquote{на підході}, --- це
поле, випромінене в момент $T-\tau$. Розкладемо поле $\Efield$ у зламі на \emphx{радіальну} (поздовжню) складову $E_\parallel$ і
\emphx{поперечну} складову $E_\perp$.

\begin{figure}[h!]\centering
	\localinput{elementary-model-derivation.tikz}
	\caption{Геометрія зламу силової лінії: розкладання поля на поздовжню $E_\parallel$ і
		поперечну $E_\perp$ складові\label{pic:ElementaryDerivation}}
\end{figure}

Із подібності трикутників (побудованих на бічних сторонах зламу) випливає:
\begin{equation}\label{eq:RatioEE}
	\frac{E_\perp}{E_\parallel} = \frac{v_qT\sin\theta}{c\tau}.
\end{equation}
Поздовжня складова --- це просто (майже незмінене) кулонівське поле заряду на відстані $r=cT$ (у гауссовій системі одиниць):
\begin{equation}\label{eq:EparDef}
	E_\parallel = \frac{q}{(cT)^2} \stackrel{r=cT}{=} \frac{q}{(cT)\,r}.
\end{equation}
Оскільки прискорення $a=v_q/\tau$, підставляючи в~\eqref{eq:RatioEE} і~\eqref{eq:EparDef}:
\begin{equation}\label{eq:EperpDerivation}
	E_\perp(t) = \frac{q}{(cT)\,r}\cdot\frac{v_qT\sin\theta}{c\tau} = \frac1{c^2}\frac{q}{r}\, a\,\sin\theta,
\end{equation}
де прискорення потрібно брати саме в той момент, коли воно й породило спостережуваний зараз злам, тобто в момент $t-r/c$:
\begin{equation}\label{eq:EperpResult}
	\tcbhighmath{
		E_\perp(t) = \frac1{c^2}\,\frac{q}{r}\, a\!\left(t-\frac{r}{c}\right)\sin\theta.
	}
\end{equation}

Формула~\eqref{eq:EperpResult} містить дві ключові риси \emphx{будь-якого} випромінювання, які ми надалі зустрічатимемо знову й знову:
\begin{enumerate}
	\item Напруженість поля \emphx{хвилі} $E_\perp$ спадає як $1/r$, на відміну від електростатичного поля $E_\parallel\propto1/r^2$. Це
	      узгоджується із законом збереження енергії: енергія хвилі розподіляється по поверхні сфери ($\propto r^2$), а густина енергії
	      $\propto E^2$, тому $E^2\cdot r^2 = \const \Rightarrow E\propto1/r$.
	\item $E_\perp$ у момент часу $t$ визначається прискоренням заряду $a$ не в цей самий момент, а в \emphx{більш ранній} момент $t-r/c$
	      --- \emphx{час запізнення}, потрібний хвилі, щоб дістатися точки спостереження на відстані~$r$.
\end{enumerate}

\subsection*{Випромінювання осцилюючого заряду}

Нехай тепер заряд $q$ не просто короткочасно прискорюється, а весь час рухається вздовж прямої зі змінним прискоренням $a(t)$ поблизу початку
координат --- наприклад, здійснює гармонічні коливання. Тоді він випромінює електромагнітні хвилі \emphx{неперервно}, і напруженість поля
випромінювання в точці на відстані $r$ від початку координат, як і раніше, визначається формулою~\eqref{eq:EperpResult}, але тепер прискорення
потрібно брати в більш ранній момент $t-r/c$:
\begin{equation}\label{eq:OscillatingChargeE}
	E_\perp(t) = \frac{1}{c^2}\,\frac{q}{r}\,a\!\left(t-\frac{r}{c}\right)\sin\theta.
\end{equation}

Нехай, зокрема, заряд здійснює гармонічні коливання з амплітудою $A$ і частотою $\omega$:
\begin{equation*}
	x(t) = A\cos\omega t, \qquad a(t) = \ddot x(t) = -\omega^2 A\cos\omega t.
\end{equation*}
Підставляючи це прискорення у~\eqref{eq:OscillatingChargeE}, отримуємо
\begin{equation}\label{eq:OscillatingChargeField}
	\tcbhighmath{
		E_\perp(t) = -\frac{q\omega^2 A}{c^2 r}\,\cos\omega\!\left(t-\frac{r}{c}\right)\sin\theta.
	}
\end{equation}
Отже, поле в кожній точці простору гармонічно коливається з тією самою частотою $\omega$: \emphx{осцилюючий заряд випромінює монохроматичну
хвилю}. Формула~\eqref{eq:OscillatingChargeField} справедлива на відстанях $r$, великих порівняно з амплітудою коливань $A$ (і, звісно, за
умови $v\ll c$).

\begin{Attention}
	\emphx{Монопольного випромінювання не існує!} Ізольований заряд не може \enquote{пульсувати} (породжувати хвилю), бо це порушило б закон
	збереження заряду. Але тоді виникає питання: як же тоді випромінює \emphx{один} заряд, що коливається? Відповідь: коли заряд коливається біля
	певного \enquote{центра}, він фактично утворює \emphx{диполь} із цим центром --- і саме дипольний момент системи й визначає випромінювання.
	Тому найпростішою системою, здатною ефективно випромінювати, є \emphx{електричний диполь} зі змінним у часі дипольним моментом. До нього ми
	й переходимо.
\end{Attention}


%% --------------------------------------------------------
\section{Дипольне випромінювання}
%% --------------------------------------------------------

\subsection*{Диполь як елементарний випромінювач}

Розглянемо електронейтральну систему з двох зарядів $+q$ і $-q$, з яких додатний --- нескінченно важкий (і тому нерухомий), а від'ємний вільно
рухається (рис.~\ref{pic:DipoleTwoCharges}). Дипольний момент такої системи та його друга похідна за часом:
\begin{equation*}
	\vect p = -q\vect r', \qquad \ddot{\vect p} = -q\dot{\vect v},
\end{equation*}
де $\vect v=\dot{\vect r}'$ --- швидкість заряду, а $\dot{\vect v}$ --- його прискорення.

\begin{figure}[h!]\centering
	\localinput{dipole-two-charges.tikz}
	\caption{Найпростіша модель диполя: важкий нерухомий заряд $+q$ і легкий рухомий $-q$\label{pic:DipoleTwoCharges}}
\end{figure}

\subsection*{Про запізнюючі потенціали (коротко)}

Строге виведення полів диполя спирається на розв'язання рівнянь Максвелла через електромагнітні потенціали $\phi$, $\vect A$ у калібруванні
Лоренца, що приводить до неоднорідного хвильового рівняння Даламбера з розв'язками у вигляді \emphx{запізнюючих потенціалів}:
\begin{equation*}
	\vect A(\vect r, t)= \frac1c \iiint\limits_{V'} \frac{\vect j\!\left(\vect r', t - \dfrac{|\vect r-\vect r'|}{c}\right) dV'}{|\vect r-\vect
	r'|}, \qquad \phi(\vect r, t) = \iiint\limits_{V'} \frac{\rho\!\left(\vect r', t - \dfrac{|\vect r-\vect r'|}{c}\right) dV'}{|\vect r-\vect
	r'|}.
\end{equation*}

Ці потенціали називають \emphx{запізнюючими}, бо поле в точці $\vect r$ у момент $t$ визначається не миттєвим розподілом зарядів і струмів, а
їхнім станом у \emphx{попередній} момент часу $t-|\vect r-\vect r'|/c$ --- з урахуванням скінченної швидкості поширення поля. Формально рівняння
Даламбера допускають і другий розв'язок, зі знаком \enquote{$+$} в аргументі (\emphx{випереджувальні потенціали}), --- але він не має ясного
фізичного сенсу (поле випереджало б причину, яка його породжує) і відкидається на основі \emphx{причинності}. Розкладаючи ці інтеграли в ряд для
випадку $r\gg r'$ (система набагато менша за відстань до спостерігача) і нерелятивістських швидкостей зарядів ($v\ll c$), можна довести, що
головний член розкладу визначається виключно \emphx{дипольним моментом} системи $\vect p(t-r/c)$ --- це й строго обґрунтовує, чому диполь є
найпростішим (і найважливішим) елементарним випромінювачем; наступні члени розкладу відповідають магнітно-дипольному, електрично-квадрупольному
й вищим типам випромінювання, зазвичай значно слабшим за дипольне.

\subsection*{Поля диполя: ближня й хвильова зони}

Нехай дипольний момент змінюється гармонічно, $\vect p = \vect p_0\cos\omega t$. Точне вираження полів (у сферичних координатах, з віссю $z$
уздовж $\vect p_0$) містить доданки різного порядку за $1/r$, і його зручно аналізувати окремо у двох граничних випадках.

\paragraph{Ближня зона.} Відстань від осцилятора до точки спостереження мала порівняно з довжиною хвилі:
\begin{equation}\label{eq:NearZoneCondition}
	r \ll \frac\lambda{2\pi}.
\end{equation}
Тут електричне поле в кожен момент часу $t$ практично збігається з полем \emphx{статичного} диполя, момент якого дорівнює миттєвому значенню
$p(t)$:
\begin{equation}\label{eq:NearZoneE}
	E_r = \frac{2\cos\theta\, p(t)}{r^3}, \qquad E_\theta = \frac{\sin\theta\, p(t)}{r^3},
\end{equation}
а магнітне поле, оскільки $\dot p(t) = Il$ (де $I$ --- еквівалентний струм, $l$ --- довжина еквівалентного елемента струму), збігається з
квазістатичним полем цього елемента струму за законом Біо--Савара--Лапласа:
\begin{equation}\label{eq:NearZoneB}
	\Bfield = \frac1{cr^3}\left[\dot{\vect p}(t)\times\vect r\right].
\end{equation}
У ближній зоні поля \emphx{не запізнюються} (запізнення $r/c \ll T = 2\pi/\omega$ --- мізерне порівняно з періодом коливань) і практично
збігаються з полями відповідного статичного диполя й елемента струму --- випромінювання тут ще \enquote{не встигло сформуватися}.

\paragraph{Хвильова зона.} Відстань, набагато більша за довжину хвилі:
\begin{equation}\label{eq:WaveZoneCondition}
	r \gg \frac\lambda{2\pi}.
\end{equation}
Тут (рис.~\ref{pic:WaveZone}) залишаються лише члени, що спадають як $1/r$ (решта --- швидше, і на великих відстанях нехтовно малі):
\begin{equation}\label{eq:WaveZoneFields}
	\tcbhighmath{
		E_\theta = B_\varphi = \frac{\sin\theta}{c^2r}\,\ddot p\!\left(t-\frac{r}c\right), \qquad E_r=E_\varphi=B_r=B_\theta=0.
	}
\end{equation}

\begin{SCfigure}[1][h!]\centering
	\includegraphics[width=0.45\linewidth]{wave_zone}
	\caption{Поля $E_\theta$ і $B_\varphi$ у хвильовій зоні диполя\label{pic:WaveZone}}
\end{SCfigure}

У хвильовій зоні поля $\Efield$ і $\Bfield$ поперечні, чисельно рівні одне одному й спадають як $1/r$ --- точнісінько ті самі риси плоскої
електромагнітної хвилі (гл.~\ref{EMWaves}). Крім того, поле залежить від полярного кута $\theta$: на осі диполя ($\theta=0,\pi$) поле
дорівнює нулю, максимуму ж воно сягає в екваторіальній площині ($\theta=\pi/2$) --- \emphx{диполь не випромінює вздовж власної осі}. У кожній
точці хвильової зони вектори $\Efield$, $\Bfield$ і $\vect r$ взаємно перпендикулярні й утворюють правовинтову трійку.

Формула~\eqref{eq:WaveZoneFields} --- \emphx{загальна}: вона описує випромінювання будь-якого диполя з довільним законом зміни моменту $p(t)$.
Зокрема, для гармонічного осцилятора $p(t)=p_0\cos\omega t$ вона дає
\begin{equation*}
	E_\theta = -\frac{p_0\omega^2}{c^2r}\,\cos\omega\!\left(t-\frac{r}{c}\right)\sin\theta,
\end{equation*}
що з точністю до позначень ($p_0=qA$) збігається з формулою~\eqref{eq:OscillatingChargeField}, отриманою вище для осцилюючого заряду. Отже,
\emphx{осцилюючий заряд є частковим випадком диполя} --- і загальна теорія дипольного випромінювання природно узагальнює результати
попереднього розділу.

\begin{SCfigure}[1][h!]
	\centering
	\includegraphics[width=0.5\linewidth]{dipole_radiation}
	\caption{Електричне поле осцилюючого диполя: у ближній зоні поле ніби \enquote{причеплене} до диполя, у
		хвильовій --- випромінюється\label{pic:DipoleRadiationSnapshot}}
\end{SCfigure}

\subsection*{Діаграма спрямованості}

Кутову залежність інтенсивності випромінювання диполя ($\propto\sin^2\theta$) зручно зображати \emphx{полярною діаграмою спрямованості}
(рис.~\ref{pic:PolarPattern}): довжина відрізка від початку координат до кривої в напрямку $\theta$ дорівнює величині поля (чи інтенсивності) у
цьому напрямку.

\begin{SCfigure}[2][h!]
	\centering
		\localinput{polar-pattern.tikz}
	\caption{Полярна діаграма спрямованості дипольного випромінювача: характерна форма
		\enquote{бублика}\label{pic:PolarPattern}}
\end{SCfigure}

\subsection*{Поляризація випромінювання}

З формули~\eqref{eq:WaveZoneFields} видно, що в хвильовій зоні відмінна від нуля лише \emphx{одна} компонента електричного поля --- $E_\theta$: вектор
$\Efield$ у кожній точці простору лежить у площині, що проходить через вісь диполя (напрямок $\vect p$) і напрямок спостереження $\vect r$. Оскільки
ця площина для заданого напрямку спостереження не змінюється з часом (диполь коливається вздовж фіксованої осі), випромінена хвиля виявляється
\emphx{лінійно поляризованою}: вектор $\Efield$ у кожній точці простору коливається вздовж однієї й тієї самої прямої, а не обертається й не описує
еліпс. Це не просто математична деталь: саме завдяки лінійній поляризації дипольного випромінювання приймальні антени ефективно вловлюють сигнал
лише тоді, коли їх зорієнтовано \emphx{вздовж} площини коливань поля хвилі (детальніше --- нижче), а поляризаційні світлофільтри здатні пропускати
чи затримувати світло певної поляризації (детальніше про поляризацію --- гл.~\ref{EMWaves}).


%% --------------------------------------------------------
\section{Потужність випромінювання осцилятора}
%% --------------------------------------------------------

Густина потоку електромагнітної енергії характеризується вектором Пойнтінга $\vect\Pi$ (гл.~\ref{MaxwellEqns}); потік енергії крізь поверхню
сфери радіуса $r$, що оточує осцилятор, і називають \emphx{потужністю випромінювання}. Елемент площі сфери в сферичних координатах $d\sigma =
r^2\sin\theta\, d\theta\, d\varphi$, тому, підставляючи поля хвильової зони~\eqref{eq:WaveZoneFields} у $\Pi = \dfrac{c}{4\pi}EB = \dfrac{c}{4\pi}E^2$:
\begin{multline}\label{eq:PowerInstant}
	\frac{dW}{dt} = \oiint\limits_\sigma \Pi\, d\sigma = \frac{c}{4\pi}\int\limits_0^{2\pi}\!\!\int\limits_0^\pi \left[\frac{\sin\theta}{c^2r}\,
	\ddot p\!\left(t-\frac{r}{c}\right)\right]^2 r^2\sin\theta\, d\theta\, d\varphi \\=\\ \frac{\ddot p^2\!\left(t-\dfrac{r}{c}\right)}{4\pi
	c^3}\cdot2\pi\int\limits_0^\pi\sin^3\theta\, d\theta.
\end{multline}
Оскільки $\displaystyle\int_0^\pi\sin^3\theta\,d\theta = \frac43$, отримуємо
\begin{equation}\label{eq:PowerInstantResult}
	\tcbhighmath{
		\frac{dW}{dt} = \frac{2\ddot p^2\!\left(t-\dfrac{r}{c}\right)}{3c^3}.
	}
\end{equation}
Тут $\ddot p$ --- друга похідна дипольного моменту за часом (тобто «прискорення» дипольного моменту). Для гармонічного осцилятора
$p=p_0\cos\omega t$, $\ddot p = -\omega^2p_0\cos\omega t$, і, усереднюючи за період ($\langle\cos^2\rangle=1/2$):
\begin{equation}\label{eq:PowerAveraged}
	\tcbhighmath{
		\left\langle\frac{dW}{dt}\right\rangle = \frac{\omega^4p_0^2}{3c^3}.
	}
\end{equation}

\begin{wrapstuff}[type=figure, o, top=2, width=0.4\linewidth]
    \centering
    \includegraphics[width=\linewidth]{HalfWaveDipole}
	\caption{Реальна півхвильова дипольна антена\label{pic:HalfWaveDipole}}
\end{wrapstuff}
Середня потужність випромінювання осцилятора пропорційна квадрату амплітуди дипольного моменту $p_0$ і --- значно різкіше --- \emphx{четвертій
степені частоти} $\omega$. Саме тому для радіозв'язку використовують відносно \emphx{короткі} хвилі (десятки метрів --- кілометри, тобто порівняно
високі частоти за мірками радіодіапазону): довгі хвилі (наднизькі частоти) від повільно змінних струмів практично не випромінюються --- $\omega^4$
у чисельнику дає надзвичайно різкий виграш у потужності випромінювання. Реальні радіоантени часто мають розмір порядку половини довжини хвилі ---
саме такі \emphx{півхвильові диполі} (рис.~\ref{pic:HalfWaveDipole}) найефективніше випромінюють і приймають електромагнітні хвилі.


\subsection*{Розсіювання світла: чому небо блакитне}


Тим же характером залежності випромінювання осцилятора від частоти пояснюється, зокрема, і блакитний колір неба (рис.~\ref{pic:BlueSky}).
Сонячне світло, що пронизує атмосферу, розсіюється молекулами повітря, які можна вважати елементарними осциляторами: під впливом світлової хвилі
вони здійснюють \emphx{вимушені} коливання. Оскільки період власних коливань цих молекул істотно відрізняється від періоду видимого світла
(резонанс відсутній), амплітуда вимушених коливань слабко залежить від частоти хвилі, тому інтенсивність розсіяного світла, тобто інтенсивність
випромінювання цих вимушених осциляторів, за формулою~\eqref{eq:PowerAveraged}, обернено пропорційна $\lambda^4$ (позаяк $\omega\propto1/\lambda$).
Тому короткохвильове (синє) світло розсіюється значно сильніше, ніж довгохвильове (червоне), --- і небо здається блакитним.

\begin{wrapstuff}[type=figure, o, top=2, width=0.5\linewidth]
    \centering
    
\includegraphics[width=\linewidth]{BlueSky}
	\caption{Розсіювання сонячного світла молекулами атмосфери --- причина блакитного кольору
		неба\label{pic:BlueSky}}
\end{wrapstuff}
Та сама залежність $\propto1/\lambda^4$ пояснює і, на перший погляд, протилежне явище --- \emphx{червоне Сонце на заході}. Коли Сонце стоїть низько
над горизонтом, його промені, перш ніж дістатися ока спостерігача, проходять крізь \emphx{набагато товщий} шар атмосфери, ніж опівдні. Синя й
фіолетова складові світла, розсіюючись найсильніше, встигають майже повністю \enquote{вибути} з прямого пучка вздовж цього довгого шляху
(розсіюючись убік --- це і є та сама блакить, яку ми бачимо, дивлячись не на Сонце, а в інший бік неба); до ока ж у прямому напрямку доходить
переважно те, що лишилося в пучку, --- довгохвильове, червоно-жовтогаряче світло. Отже, блакитне небо вдень і червоне Сонце на заході --- два
прояви \emphx{однієї й тієї самої} фізики релеївського розсіювання: розсіяне вбік світло --- блакитне, послаблений розсіюванням прямий пучок ---
червоний.

\subsection*{Опір випромінювання}

З погляду генератора, що живить антену змінним струмом, випромінена антеною енергія --- це енергія, яка \emphx{безповоротно} йде з кола, точнісінько
так само, як і енергія, що розсіюється на звичайному омічному опорі. Тому в теорії кіл зручно формально ввести \emphx{опір випромінювання}
$R_\text{вип}$ --- такий еквівалентний опір, що виділяв би в собі ту саму середню потужність $\langle dW/dt\rangle$ при тому самому амплітудному
струмі $I_0$ в антені, який реально в ній тече:
\begin{equation*}
	\left\langle\frac{dW}{dt}\right\rangle \equiv \frac12 I_0^2 R_\text{вип}.
\end{equation*}
Це дозволяє формально врахувати випромінювання нарівні зі звичайними тепловими втратами при розрахунку електричних кіл антен, хоча фізична природа
цих \enquote{втрат} зовсім інша: енергія не переходить у тепло всередині самої антени, а назавжди йде з нею у вигляді електромагнітної хвилі. Той
самий півхвильовий диполь чудово працює й у зворотному режимі --- як приймальна антена (це --- наслідок \emphx{теореми взаємності}: та сама
геометрія провідника, що ефективно перетворює змінний струм на хвилю, так само ефективно перетворює енергію набігаючої хвилі назад у струм в
антені). Саме тому характеристики антени --- діаграма спрямованості й поляризаційні властивості --- однакові й для передавання, і для приймання.


%% --------------------------------------------------------
\section{Випромінювання рухомого заряду}
%% --------------------------------------------------------

Скориставшись моделлю \enquote{важкий $+q$ --- легкий рухомий $-q$} (рис.~\ref{pic:DipoleTwoCharges}) і формулою~\eqref{eq:PowerInstantResult},
підставимо $\ddot{\vect p} = -q\dot{\vect v}$:
\begin{equation}\label{eq:LarmorFormula}
	\tcbhighmath{
			\frac{dW}{dt} = \frac{2q^2}{3c^3}\dot{\vect v}^2.
	}
\end{equation}
Тут $\dot{\vect v}$ --- прискорення заряду (друга похідна його радіус-вектора за часом). Формулу~\eqref{eq:LarmorFormula} називають
\emphx{формулою Лармора}.

\begin{Attention}
	Формула~\eqref{eq:LarmorFormula} --- один із найважливіших результатів класичної електродинаміки: \emphx{будь-який}
	заряд, що рухається з прискоренням, неминуче випромінює електромагнітну енергію, причому втрачена потужність залежить \emphx{лише} від
	заряду й миттєвого прискорення --- і зовсім не залежить від того, чому саме заряд прискорюється (електричне поле, магнітне поле, механічний
	поштовх). Заряд, що рухається \emphx{рівномірно} й прямолінійно ($\dot{\vect v}=0$), не випромінює взагалі --- випромінює лише
	\emphx{прискорений} заряд.
\end{Attention}

\subsection*{Гальмівне випромінювання}

Різке гальмування швидких заряджених частинок --- поширений і технологічно важливий прояв формули Лармора, який називають \emphx{гальмівним
випромінюванням} (нім. \textit{Bremsstrahlung}). Саме на ньому ґрунтується робота рентгенівської трубки: пучок електронів, розігнаний високою
напругою (десятки кіловольт), влучає в металевий анод і різко (за характерний час порядку розміру атома) гальмується електричним полем атомних
ядер речовини анода. За формулою~\eqref{eq:LarmorFormula} таке величезне (від'ємне) прискорення неминуче супроводжується випромінюванням; оскільки
різні електрони пучка гальмуються по-різному --- одні зупиняються практично одразу в одному зіткненні, інші --- поступово, за кілька послідовних
зіткнень, віддаючи енергію частинами, --- випромінювання виходить не на одній частоті, а суцільним спектром. Саме так рентгенівські трубки й
отримують потрібне для медичної діагностики, дефектоскопії металів та рентгеноструктурного аналізу випромінювання.

\subsection*{Синхротронне випромінювання}

Той самий принцип --- прискорено рухомий по коловій траєкторії заряд неминуче випромінює --- працює не лише для кулонівського притягання, а й
для руху заряджених частинок по колу в магнітному полі (де прискорення забезпечує сила Лоренца, а не кулонівська сила). На цьому явищі, яке
називають \emphx{синхротронним випромінюванням}, ґрунтується робота потужних лабораторних джерел рентгенівського й ультрафіолетового
випромінювання (синхротронів) для дослідження речовини й матеріалів; водночас воно ж --- головне джерело небажаних втрат енергії в кільцевих
прискорювачах заряджених частинок. Оскільки втрати різко зростають зі зменшенням маси частинки, для легких, ультрарелятивістських електронів
синхротронні втрати на коловій орбіті стають катастрофічними --- саме тому найбільші сучасні прискорювачі електронів і позитронів роблять
\emphx{лінійними}, тоді як важчі протони (Великий адронний колайдер) і далі успішно прискорюють по кільцю.

\subsection*{Стабільність атома: класична суперечність}

Формула Лармора має й драматичний історичний наслідок: електрон, що обертається по коловій орбіті навколо ядра (як у класичній
\enquote{планетарній} моделі атома Резерфорда), рухається з доцентровим прискоренням, а отже, за~\eqref{eq:LarmorFormula}, \emphx{неминуче
випромінює} й безперервно втрачає енергію. Оцінка (аналогічна до задачі~3 нижче) показує, що класичний електрон мав би \emphx{впасти на ядро}
за час порядку $10^{-11}$~с --- атоми в такій моделі просто не могли б бути стабільними!

\begin{Attention}
	Суперечність між класичною електродинамікою й спостережуваною стабільністю атомів стала одним з головних поштовхів до створення квантової
	механіки: М.~Бор 1913~р. був змушений постулювати, що на деяких (\enquote{стаціонарних}) орбітах електрон \emphx{всупереч} класичній
	електродинаміці не випромінює --- цей постулат, який не мав жодного пояснення в межах класичної фізики, згодом отримав строге обґрунтування
	вже в повній квантовій теорії.
\end{Attention}


%% --------------------------------------------------------
\section*{Задачі}
%% --------------------------------------------------------



\begin{problem}
Електромагнітна хвиля, що випромінюється диполем, поширюється у вакуумі так, що у хвильовій зоні на промені, перпендикулярному до
осі диполя, на відстані $r$ від нього середнє значення густини потоку енергії дорівнює $\langle\Pi\rangle$. Знайти середню потужність
випромінювання диполя.
\end{problem}

\paragraph{Розв'язання.} У хвильовій зоні на промені, перпендикулярному до осі диполя ($\theta=\pi/2$), $\sin\theta=1$, і
\begin{equation*}
	\langle\Pi\rangle = \frac{c}{4\pi}\langle E^2\rangle = \frac{c}{4\pi}\cdot\frac{\langle\ddot p^2\rangle}{c^4r^2}
	= \frac{\langle\ddot p^2\rangle}{4\pi c^3r^2}.
\end{equation*}
Звідси $\langle\ddot p^2\rangle = 4\pi c^3r^2\langle\Pi\rangle$. Підставляючи в~\eqref{eq:PowerInstantResult}, усереднене за часом:
\begin{equation*}
	\tcbhighmath{
		\left\langle\frac{dW}{dt}\right\rangle = \frac{2}{3c^3}\cdot 4\pi c^3r^2\langle\Pi\rangle = \frac{8\pi r^2}{3}\langle\Pi\rangle.
	}
\end{equation*}


\begin{problem}
Сталий за модулем електричний дипольний момент $\vect p$ обертають з кутовою швидкістю $\omega$ навколо осі, перпендикулярної до
осі диполя й що проходить через його середину. Знайти потужність випромінювання диполя.
\end{problem}

\paragraph{Розв'язання.} Розкладемо вектор диполя на дві взаємно перпендикулярні складові в площині обертання: $p_x=p\cos\omega t$,
$p_y=p\sin\omega t$. Тоді
\begin{multline*}
	\ddot p_x = -p\omega^2\cos\omega t, \qquad \ddot p_y = -p\omega^2\sin\omega t, \\ \ddot{\vect p}^{\,2} = \ddot p_x^2+\ddot p_y^2 =
	p^2\omega^4\left(\cos^2\omega t+\sin^2\omega t\right) = p^2\omega^4 = \const.
\end{multline*}
За формулою~\eqref{eq:PowerInstantResult} (яка тут одразу постійна в часі, тож усереднення не потрібне):
\begin{equation*}
	\tcbhighmath{
		\frac{dW}{dt} = \frac{2p^2\omega^4}{3c^3}
	}
\end{equation*}
--- рівно вдвічі більше, ніж у лінійного осцилятора~\eqref{eq:PowerAveraged} з тією самою амплітудою $p$.

\paragraph{Зауваження.} Обертовий диполь --- це суперпозиція \emphx{двох} лінійних осциляторів у взаємно перпендикулярних напрямках, зсунутих за
фазою на $90^\circ$. Оскільки миттєва потужність кожного окремого осцилятора коливається в протифазі до іншої (коли $\cos^2\omega t$ максимальний,
$\sin^2\omega t$ --- мінімальний, і навпаки), їхня \emphx{сума} виявляється сталою в часі, а не лише в середньому, --- обертовий диполь
випромінює цілком рівномірно, без пульсацій потужності.


\begin{problem}
Знайти потужність випромінювання нерелятивістської частинки із зарядом $e$ і масою $m$, що рухається коловою орбітою радіуса $R$
у полі нерухомого точкового заряду $q$.
\end{problem}

\paragraph{Розв'язання.} Кулонівська сила забезпечує доцентрове прискорення: $\dfrac{eq}{R^2} = \dfrac{mv^2}{R}$, звідки доцентрове прискорення
\begin{equation*}
	\dot v = \frac{v^2}{R} = \frac{eq}{mR^2}.
\end{equation*}
Підставляючи у формулу Лармора~\eqref{eq:LarmorFormula}:
\begin{equation*}
	\tcbhighmath{
		\frac{dW}{dt} = \frac{2e^2}{3c^3}\left(\frac{eq}{mR^2}\right)^2 = \frac{2e^4q^2}{3c^3m^2R^4}.
	}
\end{equation*}

\paragraph{Зауваження.} Саме ця формула (з $e=q=$ заряд електрона, $m=$ маса електрона, $R\sim10^{-10}$~м --- розмір атома) і дає ту саму
драматичну оцінку часу падіння класичного електрона на ядро $\sim10^{-11}$~с, про яку йшлося вище. Той самий принцип працює й для руху заряджених
частинок по колу в магнітному полі (сила Лоренца замість кулонівської) --- це \emphx{синхротронне випромінювання}, про яке йшлося вище.


%% --------------------------------------------------------
\section*{Підсумки}
%% --------------------------------------------------------

%\begin{table}[h!]\small\centering
%	\caption{Основні співвідношення теорії випромінювання}
%	\begin{tabular}{ll}
%		\toprule
%		Поле хвильової зони диполя      & $\displaystyle E_\theta = B_\varphi = \frac{\sin\theta}{c^2r}\,\ddot p\!\left(t-\frac rc\right)$ \\
%		Миттєва потужність випромінювання диполя & $\displaystyle \frac{dW}{dt} = \frac{2\ddot p^2}{3c^3}$                                 \\
%		Середня потужність гармонічного осцилятора & $\displaystyle \left\langle\frac{dW}{dt}\right\rangle = \frac{\omega^4p_0^2}{3c^3}$     \\
%		Формула Лармора (довільний заряд) & $\displaystyle \frac{dW}{dt} = \frac{2q^2}{3c^3}\dot{\vect v}^2$                                \\
%		\bottomrule
%	\end{tabular}
%\end{table}

Розділ про випромінювання замикає весь курс електрики й магнетизму найглибшим можливим підсумком: \emphx{будь-який} прискорений заряд
неминуче породжує хвилю, що назавжди йде від нього в навколишній простір, --- а електромагнітне поле остаточно постає перед нами не як
допоміжна математична конструкція для опису сил між зарядами, а як цілком самостійна, фізично реальна субстанція, здатна існувати,
поширюватись і переносити енергію \emphx{без} будь-яких зарядів поблизу. Саме на цьому явищі --- випромінюванні прискорених зарядів ---
ґрунтується вся радіотехніка, вся астрофізика випромінювання зір і туманностей, і, зрештою, саме воно --- а не якийсь окремий, спеціально
влаштований \enquote{носій світла} --- відповідає за існування й поширення видимого світла в природі.